% from ega4-2-fr.tex
\oldpage[IV-2]{185}
En effet, $\widehat M=M\otimes_A\widehat A$ est un $\widehat A$-module de
Cohen--Macaulay (0, 16.5.2), de support égal à
$\operatorname{Spec}(\widehat A)$; par suite (0, 16.5.4), $\widehat A$ est
équidimensionnel.
\end{proof}

\begin{remark}[7.1.6]
\label{IV.2.s7.1.6-fr}
On a vu (6.3.8) que sous les hypothèses de (7.1.5), pour tout anneau quotient
$B$ de $A$, les fibres du morphisme canonique
$\operatorname{Spec}(\widehat B)\to\operatorname{Spec}(B)$ sont des schémas
de Cohen--Macaulay; par suite ((6.4.3) et (5.7.5)), \emph{pour que $B$ n'ait
pas d'idéaux premiers associés immergés, il faut et il suffit qu'il en soit
de même de $\widehat B$}.
\end{remark}

\begin{lemma}[7.1.7]
\label{IV.2.s7.1.7-fr}
Soient $A$ un anneau local noethérien, $\mathfrak p_i$
($1\leqslant i\leqslant r$) ses idéaux premiers minimaux. Supposons que
chacun des anneaux $A/\mathfrak p_i$ soit formellement équidimensionnel.
Alors, pour tout idéal premier $\mathfrak q$ de $A$ et tout $i$ tel que
$\mathfrak p_i\subset\mathfrak q$, l'anneau
$A_{\mathfrak q}/\mathfrak p_iA_{\mathfrak q}$ est formellement
équidimensionnel.
\end{lemma}

\begin{proof}
Comme $A_{\mathfrak q}/\mathfrak p_iA_{\mathfrak q}$ est l'anneau local de
$A/\mathfrak p_i$ en l'idéal premier $\mathfrak q/\mathfrak p_i$, on peut
se borner au cas où $A$ est intègre et formellement équidimensionnel. Posons
$A'=\widehat A$, et soit $\mathfrak q'$ un des idéaux premiers de $A'$
minimaux parmi ceux contenant $\mathfrak qA'$; si l'on pose
$B=A_{\mathfrak q}$, $B'=A'_{\mathfrak q'}$, $B'$ est un $B$-module plat
(0$_{\mathrm I}$, 6.3.2). Posons $C=\widehat B$, $C'=\widehat {B'}$; comme
$B'$ est un $B$-module plat, on sait que $C'$ est un $C$-module plat
(Bourbaki, \emph{Alg. comm.}, chap. III, § 5, n\textsuperscript{o} 4,
prop. 4). Comme $A'$ est caténaire (5.6.4) et équidimensionnel par
hypothèse, il en est de même de $B'=A'_{\mathfrak q'}$ (0, 16.1.4); en
outre, comme $A'$ est isomorphe à un quotient d'anneau régulier en vertu du
théorème de Cohen (0, 19.8.8), il en est de même de $B'$ (0, 17.3.9); on
conclut donc de (7.1.5) que $C'$ est équidimensionnel. D'autre part $C'$ est
caténaire (5.6.4), donc $C$ est équidimensionnel en vertu de (7.1.3, (ii)).
\end{proof}

\begin{theorem}[7.1.8]
\label{IV.2.s7.1.8-fr}
Soit $A$ un anneau local noethérien. Les conditions suivantes sont
équivalentes:
\begin{enumerate}
\item[a)] Tout anneau quotient intègre de $A$ est formellement équidimensionnel.
\item[b)] Les anneaux quotients de $A$ par ses idéaux premiers minimaux sont
formellement équidimensionnels.
\item[c)] Tout anneau local $B$ qui est une $A$-algèbre essentiellement de type
fini (1.3.8) et qui est équidimensionnel, est formellement
équidimensionnel.
\end{enumerate}
\end{theorem}

\begin{proof}
Il est trivial que c) entraîne a) et que a) entraîne b). D'autre part,
comme tout idéal premier de $A$ contient un idéal premier minimal de $A$,
b) entraîne a) en vertu de (7.1.4, (ii)). Il reste à voir que a) implique
c).

Il suffit de prouver que les quotients de $B$ par ses idéaux premiers
minimaux sont formellement équidimensionnels (7.1.2); comme tout anneau
quotient de $B$ est encore une $A$-algèbre essentiellement de type fini
(1.3.9), on peut en premier lieu supposer $B$ \emph{intègre}. Si
$\mathfrak q$ est
l'image réciproque dans $A$ de l'idéal maximal de $B$, on sait que $B$ est
aussi une $A_{\mathfrak q}$-algèbre essentiellement de type fini (1.3.10),
et il résulte de a) et de (7.1.7) que tout anneau quotient intègre de
$A_{\mathfrak q}$ est formellement équidimensionnel; on peut donc supposer
que l'homomorphisme $A\to B$ est un homomorphisme \emph{local}. Le noyau
$\mathfrak r$ de cet homomorphisme est alors un idéal premier de $A$, et en
vertu de a), tout anneau quotient intègre de $A/\mathfrak r$ est
formellement équidimensionnel; comme $B$ est une
$(A/\mathfrak r)$-algèbre essentiellement de type fini, on voit qu'on peut
aussi supposer

